\begin{textsample}{POS Dim 5 – vem\_ed\_27 – Score 3.00 – vem\_ed\_27/t142.txt}  \label{ex:f5_pos_037}
Aos Leitores Osvaldo Succi Jr.
Coordenador PCIs Os Projetos Colaborativos Internacionais (PCI / Cesu) representam uma interpretação das Fatecs do Centro Paula Souza (CPS) de uma abordagem de Intercâmbios Virtuais conhecida internacionalmente como COIL (Collaborative Online Internacional Learning – Aprendizagem Internacional Colaborativa realizada on-line).
Resumidamente, professores de Fatecs e de \textbf{universidades} internacionais propõem aos alunos atividades pedagógicas para serem desenvolvidas em equipes mistas (brasileiros e estrangeiros).
O primeiro PCI / Cesu foi realizado em 2013 na Fatec Americana e, em 2014, venceu o Prêmio Santander Universidades de Melhor Parceria Acadêmica.
Em 2024, ao completar dez anos dessa conquista, o Centro Paula Souza é reconhecido como referência em COIL no mundo.
A Jornada de PCIs 2024, realizada em dezembro, celebra esse prestígio e compartilha boas práticas de professores de Fatecs e Instituições de Ensino Superior (IES) brasileiras, bem como reflexões e pesquisas em Intercâmbios Virtuais e políticas para Internacionalização em Casa.
O evento está detalhado na reportagem principal desta edição.
Entre 2013 e 2024, 12.630 fatecanos participaram de 630 PCIs em 53 Fatecs.
As colaborações foram realizadas com 59 IES de 18 países.
A expansão dos PCIs reflete a consolidação das ações desenvolvidas pela equipe dos PCIs / Cesu na prospecção de colaborações, orientação ao desenvolvimento de projetos, acompanhamento e avaliação dos resultados.
E muito desse sucesso vem do engajamento dos professores das Fatecs e do diferencial dos \textbf{currículos} dos Cursos Superiores de Tecnologia do CPS, com aulas de inglês e espanhol na matriz curricular, que facilitam a comunicação em equipes internacionais.
A proeminência das Fatecs do CPS no universo se reflete na visibilidade internacional alcançada em apresentações dos membros da equipe PCI / Cesu em congressos internacionais, a exemplo de Faubai e IVEC; convites para participar do conselhos executivos de \textbf{redes} internacionais (COIL Connect e Red LatAm COIL) e para ministrar cursos de formação na abordagem COIL para docentes e gestores de \textbf{universidades}.
Em dezembro de 2024, estive no México e no Chile, ministrando palestras e capacitações, resumidas nesta edição.
Boa leitura!

% matched lemmas: currículo, rede, universidade
\end{textsample}
